%% kapa: exam=1 round=32 year=2021 period=2 form=A subject=accounting subject_ko=회계학 seq=13 number=53
%% source: Q-Net 감정평가사 기출문제 (cst00309) artlSeq=5208397 file=extracted/1008/5208397/A형/2021년도 제32회 감정평가사 1차 2교시 A형.pdf
%% conversion: kordoc 4.24.51 -> tools/gen1.py (auto); review=proofread-sonnet; flags=none
\begin{question}{53}
재무제표의 표시에 관한 설명으로 옳지 않은 것은?
\begin{choices}
  \item 재무제표가 한국채택국제회계기준의 요구사항을 모두 충족한 경우가 아니라면 한국채택국제회계기준을 준수하여 작성되었다고 기재하여서는 안 된다.
  \item 기업이 재무상태표에 유동자산과 비유동자산으로 구분하여 표시하는 경우, 이연법인세자산은 유동자산으로 분류하지 아니한다.
  \item 비용을 기능별로 분류하는 기업은 감가상각비, 기타 상각비와 종업원급여비용을 포함하여 비용의 성격에 대한 추가 정보를 공시한다.
  \item 수익과 비용의 어느 항목은 포괄손익계산서 또는 주석에 특별손익항목으로 별도 표시한다.
  \item 매출채권에 대한 대손충당금을 차감하여 관련 자산을 순액으로 측정하는 것은 상계표시에 해당하지 아니한다.
\end{choices}
\end{question}
\answer{4}
